\documentclass[UTF8,fontset=fandol]{ctexart}

\usepackage{amsmath,tikz}

\begin{document}

\section*{追逐运动：把两个速度投影合成一个导数}

目标点 $F$ 水平向右以恒定速率 $v_F$ 运动，追逐点 $D$ 以恒定速率 $v_D>v_F$ 沿视线追赶。初始两点上下对齐，相距 $L$。图 1 画出追赶中的一个时刻。求怎样把夹角从间距 $r$ 和水平领先量 $x$ 的变化率中消去。

\begin{center}\begin{tikzpicture}[>=stealth]
\coordinate (D) at (0,0); \coordinate (F) at (3.5,2);
\draw[dashed] (D)--(3.5,0)--(F);
\draw[blue,thick] (D)--node[above]{$r$}(F);
\draw[->] (F)--++(1.2,0) node[right]{$v_F$};
\draw[->,thick] (D)--(1.1,.63) node[above left]{$v_D$};
\fill (D)circle(1.5pt) node[left]{$D$};\fill (F)circle(1.5pt) node[above]{$F$};
\draw (.8,0) arc (0:30:.8) node[right]{$\theta$};\node at (1.75,-.28){$x$};
\end{tikzpicture}
\\ 图 1\end{center}

沿图 1 的视线投影，目标点使间距增加，追逐点使间距减小。

\begin{equation}\dot r=v_F\cos\theta-v_D\label{eq:s1}\end{equation}

沿水平方向，蓝点的速度减去绿点的水平分量给出领先量的变化率。

\begin{equation}\dot x=v_F-v_D\cos\theta\label{eq:s2}\end{equation}

给第一式两边同乘追逐点速率。

\begin{equation}v_D\dot r=v_Dv_F\cos\theta-v_D^2\label{eq:s3}\end{equation}

给第二式两边同乘目标点速率。

\begin{equation}v_F\dot x=v_F^2-v_Dv_F\cos\theta\label{eq:s4}\end{equation}

两式相加，两个余弦项抵消。

\begin{equation}v_D\dot r+v_F\dot x=v_F^2-v_D^2\label{eq:s5}\end{equation}

追逐点速率恒定，可以移入时间导数。

\begin{equation}v_D\dot r=\frac{d}{dt}(v_Dr)\label{eq:s6}\end{equation}

用同样的方法把目标点速率移入时间导数。

\begin{equation}v_F\dot x=\frac{d}{dt}(v_Fx)\label{eq:s7}\end{equation}

两个导数相加就是和的导数。追上时 $r=x=0$，因此这条关系也给出了求追及时间的积分条件。

\begin{equation}\frac{d}{dt}(v_Dr+v_Fx)=v_F^2-v_D^2\label{eq:s8}\end{equation}

\end{document}
